\subsection{The $SO(32)$ embedding}

The $SO(32)$ gauge group of Type~I / heterotic string theory
admits the maximal-subgroup chain
$SO(32) \supset SU(16) \times U(1)_A
  \supset SU(5) \times SU(3) \times U(1)_A \times U(1)_B$,
where $\mathbf{16} = (\mathbf{5},\mathbf{3})_{+1}
  \oplus (\mathbf{1},\mathbf{1})_{-15}$
embeds the five quark flavors times three colors as a single
tensor-product block.

\paragraph{Complete decomposition.}
The adjoint $\mathbf{496} = \wedge^2(\mathbf{32})$ splits as
$(\mathbf{255}\!\oplus\!\mathbf{1})_{0}
  \oplus \mathbf{120}_{+2}
  \oplus \overline{\mathbf{120}}_{-2}$
under $SU(16) \times U(1)_A$.  The $SU(5) \times SU(3)$ content:

\begin{center}
\begin{tabular}{llcccl}
\toprule
Sector & $SU(5)\times SU(3)$ & dim & $q_A$ & $q_B$ & Origin \\
\midrule
$\mathbf{255}\!\oplus\!\mathbf{1}$
 & $(\mathbf{24},\mathbf{8})$ & 192 & 0 & 0
   & $\mathrm{adj}(SU(15))$ \\
 & $\boldsymbol{(24,1)}$ & $\boldsymbol{24}$ & 0 & 0
   & $\mathrm{adj}(SU(15))$ \\
 & $(\mathbf{1},\mathbf{8})$ & 8 & 0 & 0
   & $\mathrm{adj}(SU(15))$ \\
 & $(\mathbf{5},\mathbf{3})\!+\!(\bar{\mathbf{5}},\bar{\mathbf{3}})$
   & 30 & 0 & $\pm 16$ & cross-terms \\
 & $2\times(\mathbf{1},\mathbf{1})$ & 2 & 0 & 0
   & $U(1)_{A,B}$ \\
\midrule
$\mathbf{120}$
 & $\boldsymbol{(15,\bar{3})}$ & $\boldsymbol{45}$ & $+2$ & $+2$
   & $S^2(\mathbf{5})\otimes\wedge^2(\mathbf{3})$ \\
 & $(\overline{\mathbf{10}},\mathbf{6})$ & 60 & $+2$ & $+2$
   & $\wedge^2(\mathbf{5})\otimes S^2(\mathbf{3})$ \\
 & $(\mathbf{5},\mathbf{3})$ & 15 & $+2$ & $-14$
   & cross-term \\
\midrule
$\overline{\mathbf{120}}$
 & $\boldsymbol{(\overline{15},3)}$ & $\boldsymbol{45}$ & $-2$ & $-2$
   & conjugate \\
 & $(\mathbf{10},\bar{\mathbf{6}})$ & 60 & $-2$ & $-2$
   & conjugate \\
 & $(\bar{\mathbf{5}},\bar{\mathbf{3}})$ & 15 & $-2$ & $+14$
   & conjugate \\
\midrule
 & \textbf{Total} & \textbf{496} & & & \\
\bottomrule
\end{tabular}
\end{center}

\noindent
The target---$(\mathbf{24},\mathbf{1})$,
$(\mathbf{15},\bar{\mathbf{3}})$,
$(\overline{\mathbf{15}},\mathbf{3})$---accounts
for $24 + 45 + 45 = 114$ of the 496 states, each appearing once.

\paragraph{Why the symmetric $\mathbf{15}$ appears.}
The nontrivial content of $\wedge^2(\mathbf{16})$ comes from
$\wedge^2[(\mathbf{5},\mathbf{3})]$.  The key identity is
$$
\wedge^2(V \otimes W)
  \;=\; \bigl[S^2(V) \otimes \wedge^2(W)\bigr]
  \;\oplus\; \bigl[\wedge^2(V) \otimes S^2(W)\bigr]\,.
$$
Overall antisymmetry of the composite index $(ia)$ forces a
correlation: symmetric in one factor implies antisymmetric in the
other.  Setting $V = \mathbf{5}$, $W = \mathbf{3}$:
$$
\wedge^2\bigl[(\mathbf{5},\mathbf{3})\bigr]
  \;=\; (\mathbf{15},\bar{\mathbf{3}})
  \;\oplus\; (\overline{\mathbf{10}},\mathbf{6})\,.
$$
The symmetric flavor $\mathbf{15}$ is forced into antisymmetric
color $\bar{\mathbf{3}}$; the pairing is dictated by the
tensor-product structure, not by a modeling choice.

\paragraph{Minimality and obstructions.}
The identity requires $(\mathbf{5},\mathbf{3})$ as a
tensor-product block in the fundamental, hence dimension
$\geq 15$ complex $= 30$ real; $SO(30)$ is the minimal
orthogonal group, and $SO(32)$ adds only two inert singlets.
$SO(16) \times SO(16)$ fails: no single $\mathbf{16}$ can hold
the 15-dimensional block, so $SU(5)$ and $SU(3)$ embed as
separate direct summands whose $\wedge^2$ produces only
$\overline{\mathbf{10}}$ and $\bar{\mathbf{3}}$, never
$\mathbf{15}$.
$E_8 \times E_8$ is likewise blocked: every $E_8$ decomposition
under $SU(5)$ yields only exterior powers
($\mathbf{5}$, $\overline{\mathbf{10}}$, $\mathbf{24}$, \ldots).
The symmetric $\mathbf{15} = S^2(\mathbf{5})$ is algebraically
forbidden.

\paragraph{Orientifold projection.}
Wilson lines and orbifold projections cannot separate
$(\mathbf{15},\bar{\mathbf{3}})$ from
$(\overline{\mathbf{10}},\mathbf{6})$: both carry identical
$U(1)$ charges $(q_A, q_B) = (+2,+2)$, and Schur's lemma forces
any orbifold element commuting with $SU(5) \times SU(3)$ to act
as a scalar on $(\mathbf{5},\mathbf{3})$ and hence on its exterior
power.
A mixed orientifold resolves this.  Placing the $SU(5)$ flavor
branes near an $O^-$ plane (Sp-type, selecting symmetric
$\mathbf{15}$) and the $SU(3)$ color branes near an $O^+$ plane
(SO-type, selecting antisymmetric $\bar{\mathbf{3}}$), worldsheet
parity~$\Omega$ acts with opposite signs on the two tensor
factors and doubly projects out
$(\overline{\mathbf{10}},\mathbf{6})$.  The remaining
$(\mathbf{24},\mathbf{8})$ and bifundamental copies acquire
compactification-scale masses from brane separation; the $U(1)$
generators become massive via the St\"uckelberg mechanism.

\paragraph{Open problem.}
An explicit compactification with the mixed $O^-/O^+$ structure
and correct tadpole cancellation has not been constructed.
The self-referential numerology---$5 \times 3 + 1 = 16$
Chan-Paton factors forming $SO(32)$, whose adjoint contains
exactly the composites of those quarks---awaits a complete
string vacuum.
